\begin{table}[htbp]
\caption{Three valid certificate accounts for the same frozen fitted policies}\label{tab:accounts95}
\centering\small
\begin{tabular}{@{}llrrrr@{}}\toprule
$\alpha$ & Continuation & Supremum norm & Signed scalar & Joint support & Qualified\\\midrule
0.5 & Quadratic & 0.981430 & 0.617589 & 0.005501 & 2048/2,048\\
0.5 & ReLU & 0.162986 & 0.122750 & 0.021836 & 2048/2,048\\
0.5 & Spline & 0.943607 & 0.581660 & 0.012112 & 2048/2,048\\
0.75 & Quadratic & 1.005610 & 0.612983 & 0.006422 & 2048/2,048\\
0.75 & ReLU & 0.170090 & 0.128429 & 0.022145 & 2048/2,048\\
0.75 & Spline & 0.847773 & 0.521923 & 0.008098 & 2048/2,048\\
\bottomrule\end{tabular}
\par\smallskip\begin{minipage}{.98\linewidth}\footnotesize
Entries are maximum original-optimum policy-loss bounds across the same 2,048 rows. Supremum-norm and signed-scalar bounds use the same frozen whole-square residual cover. Joint support uses the paid original-value witness. No row satisfies $1/16$ with either scalar account; every row does with the joint certificate. This is post-evaluation deterministic re-accounting, not a prospective runtime ablation or an optimality claim for any certificate-construction algorithm.
\end{minipage}\end{table}
